\section{Conclusion}
\label{sec:concl}

Matchertext is a syntactic discipline
that enables the verbatim embedding of strings across languages
with no transformation or expansion,
by enforcing only the single rule that the ASCII matchers --
parentheses, square brackets, and curly braces --
must always appear strictly in matched pairs.
We develop the basic principles and rationale for matchertext,
and explore potential backward-compatible extensions
to existing classes of languages
both to host embedded matchertext strings,
and to write embeddable matchertext conveniently.
We give the discipline a machine-checked foundation,
proving in Lean that delimiting any valid matchertext
again yields valid matchertext (\cref{sec:design:abstract:formal}),
so embeddings compose to any depth without escaping or expansion.
Building on that structural guarantee,
we show how matcher-delimiting untrusted input
turns defenses against injection attacks --
SQL injection, cross-site scripting, and their LDAP and PDF analogues --
into a single host-independent check (\cref{sec:host:inject}).
An empirical analysis of roughly 470~million string literals
from thousands of production repositories (\cref{sec:eval})
finds real-world code already overwhelmingly compliant with the discipline,
and that adopting it would remove around 95\% of the backslash escapes
in string literals,
evidence that the adoption cost is small while the benefit is large.
The chief remaining question is one of usability:
how matchertext affects developers writing and editing
common embedded constructs in practice,
which we leave to future work.

\subsection*{Acknowledgments}

I wish to thank
Antoine Bastide,
Henry Corrigan-Gibbs,
Russ Cox,
Philip Hamelink,
and
Terence Parr
for valuable feedback on earlier drafts of this paper.

\xxx{samller contributions and funding acknowledgments here}

